\clearpage
\fancyhead[L]{\small HARDYS PENDELUHR.}
\hypertarget{native-458-hardy}{}
\begin{center}{\large HARDYS PENDELUHR.}\end{center}

\begin{center}
\rule{0.55\linewidth}{0.4pt}\\[2pt]
Göttingische gelehrte Anzeigen. Stück 101. Seite 1001\,..\,1002. 1826. Juni 26.\\[2pt]
\rule{0.55\linewidth}{0.4pt}
\end{center}

Seine Königliche Hoheit, der Herzog von \textsc{Sussex}, hat der Universität erneut sein gnädiges Wohlwollen gezeigt, indem er der Sternwarte ein kostbares Geschenk machte. Es ist eine astronomische Pendeluhr von \textsc{Hardy}. Seit einigen Monaten dient sie bereits als Hauptuhr bei den Meridianbeobachtungen. Dabei hat sie die Erwartungen, die der Name des Uhrmachers und die Vollkommenheit seiner Arbeit geweckt hatten, voll erfüllt. Wir fügen eine Übersicht über ihren bisherigen Gang bei, dessen Gleichmäßigkeit nichts zu wünschen übrig lässt. Die Kompensation wird auf \SourceLeaf{459}denkbar einfachste Weise erreicht. Die Pendelstange trägt am unteren Ende statt der sonst üblichen Pendellinse einen Rahmen. Darin steht ein zylindrisches Glasgefäß von $6\frac{1}{2}$ Zoll Höhe und 2 Zoll Weite, das fast ganz mit Quecksilber gefüllt ist. Erweist sich die Kompensation als zu schwach, muss etwas Quecksilber hinzugegeben werden. Im umgekehrten Fall muss man etwas herausnehmen. So lässt sich die Kompensation sehr fein einstellen, sobald der Gang der Uhr bei den äußersten Temperaturen beobachtet worden ist. Die folgende Übersicht zeigt, dass der tägliche Gang bei zunehmender Wärme ganz allmählich um eine halbe Sekunde zurückgeht und die Kompensation somit bisher noch etwas zu schwach ist.

\begin{center}
\emph{Stand und Gang des Regulators von \textup{\textsc{Hardy}} gegenüber der Sternzeit, auf den Mittag reduziert.}
\end{center}

\begin{center}
\begingroup\small\setlength{\tabcolsep}{5pt}
\begin{tabular}{lr|r|r}
 & & Stand & täglicher Gang \\
\hline
1826 & März 9 & $+0^{\mathrm s}83$ & $+0.42$ \\
 & 11 & $1.67$ & $+0.52$ \\
 & 14 & $3.22$ & $+0.45$ \\
 & 19 & $5.46$ & $+0.35$ \\
 & April 20 & $16.55$ & $+0.14$ \\
 & 21 & $16.69$ & $+0.07$ \\
 & 22 & $16.76$ & $+0.21$ \\
 & Mai 2 & $19.09$ & $+0.16$ \\
 & 8 & $20.05$ & $+0.18$ \\
 & 9 & $20.23$ & $+0.25$ \\
 & 11 & $20.74$ & $+0.01$ \\
\end{tabular}
\par\medskip
\begin{tabular}{lr|r|r}
 & & Stand & täglicher Gang \\
\hline
1826 & Mai 14 & $+20^{\mathrm s}78$ & $+0.08$ \\
 & 16 & $20.94$ & $+0.21$ \\
 & 18 & $21.37$ & $+0.15$ \\
 & 20 & $21.68$ & $+0.02$ \\
 & 23 & $21.73$ & $+0.10$ \\
 & 26 & $22.03$ & $-0.06$ \\
 & 30 & $21.78$ & $-0.18$ \\
 & Juni 11 & $19.57$ & $-0.22$ \\
 & 16 & $18.45$ & $-0.18$ \\
 & 17 & $18.27$ &  \\
\end{tabular}
\endgroup
\end{center}

